\textbf{Kurzfassung}: Neben Anwendung als wissenschaftliche
Forschungsmethode können Expert\*inneninterviews in Bibliotheken als
Möglichkeit herangezogen werden, um vorhandenes Fachwissen zu
dokumentieren, damit ein aktueller reibungsloser Arbeitsablauf im
Betrieb gewährleistet werden kann. Zudem gestattet solch eine
Dokumentation neuen Mitarbeiter\*innen, sich möglichst schnell im neuen
Wirkungsbereich zurechtzufinden. Eine weitere Anwendung finden
Interviews in Bibliotheken für den Erfahrungsaustausch unter
Fachkolleg\*innen. Die Autor\*innen testen als ehemalige direkte
Kolleg*innen letztere Variante zum Thema Neukonzeption von Bibliotheken
im Zuge eines Arbeitsstellenwechsels. Der besseren Lesbarkeit willen
erfolgen alle weiteren personenbezogenen Bezeichnungen in diesem Aufsatz
in der männlichen Form, welche aber keinerlei wertende Bedeutung hat und
alle Geschlechter gleichermaßen repräsentiert.

\begin{center}\rule{0.5\linewidth}{0.5pt}\end{center}

\noindent\textbf{Abstract}: Expert interviews in libraries can be used as a
scientific research method, but also for documenting existing expertise
for ensuring smooth and efficient workflows. Such documentation also
allows new employees to quickly familiarize themselves with their new
work environment. Another application of interviews in libraries is the
exchange of experiences among colleagues. As former direct colleagues,
the authors are testing this latter approach in the context of library
redesign following off a job change.
